\documentclass[11pt, letterpaper]{article}

\usepackage[utf8]{inputenc}
\usepackage[T1]{fontenc}
\usepackage[spanish, es-tabla]{babel}
\usepackage{lmodern}
\usepackage{geometry}
\usepackage{setspace}
\usepackage{amsmath, amssymb}
\usepackage{booktabs}
\usepackage{longtable}
\usepackage{tabularx}
\usepackage{array}
\usepackage{xcolor}
\usepackage{listings}
\usepackage{hyperref}
\usepackage{enumitem}
\usepackage{fancyhdr}

\geometry{
    left=2.2cm,
    right=2.2cm,
    top=2.4cm,
    bottom=2.4cm,
    headheight=18pt,
    headsep=0.45cm
}

\setstretch{1.18}
\setlength{\parindent}{0pt}
\setlength{\parskip}{0.45em}
\renewcommand{\arraystretch}{1.15}
\setlist[itemize]{leftmargin=1.2em}
\setlist[enumerate]{leftmargin=1.4em}

\hypersetup{
    colorlinks=true,
    linkcolor=black,
    citecolor=black,
    urlcolor=blue
}

\fancypagestyle{guia}{
    \fancyhf{}
    \fancyhead[L]{Gu\'ia de estudio para defensa}
    \fancyhead[R]{\thepage}
    \renewcommand{\headrulewidth}{0.4pt}
}

\pagestyle{guia}

\definecolor{codebg}{RGB}{248,248,248}
\definecolor{codeframe}{RGB}{220,220,220}
\definecolor{codecomment}{RGB}{80,120,80}
\definecolor{codekeyword}{RGB}{0,70,140}
\definecolor{codestring}{RGB}{150,60,30}

\lstdefinestyle{pythonstyle}{
    language=Python,
    backgroundcolor=\color{codebg},
    frame=single,
    rulecolor=\color{codeframe},
    basicstyle=\ttfamily\footnotesize,
    keywordstyle=\color{codekeyword}\bfseries,
    commentstyle=\color{codecomment},
    stringstyle=\color{codestring},
    showstringspaces=false,
    breaklines=true,
    columns=fullflexible,
    tabsize=4
}

\newcommand{\fd}{f_d}
\newcommand{\pnoct}{P_{\mathrm{NOCT}}}
\newcommand{\pref}{P_{\mathrm{ref}}}
\newcommand{\phybrid}{P_{\mathrm{hybrid}}}

\begin{document}

\begin{titlepage}
    \centering
    \vspace*{1.2cm}
    {\Large Universidad Andr\'es Bello\par}
    {\large Ingenier\'ia F\'isica\par}
    \vspace{1.5cm}
    {\LARGE \textbf{Gu\'ia de estudio para defensa}\par}
    \vspace{0.4cm}
    {\Large Proyecto de t\'itulo fotovoltaico NOCT + XGBoost\par}
    \vspace{1.2cm}
    {\large Nicol\'as Antonio Aravena Osorio\par}
    \vfill
    {\large Santiago, Chile\par}
    {\large 2026\par}
\end{titlepage}

\tableofcontents
\newpage

\section{Variables que debes tener clar\'isimas}

Esta secci\'on es para estudiar. La idea es que puedas responder r\'apido qu\'e significa cada variable, de d\'onde sale, c\'omo se usa y por qu\'e importa.

\subsection{Variables solares y meteorol\'ogicas}

\begin{longtable}{p{3.0cm}p{5.1cm}p{6.0cm}}
\toprule
\textbf{Variable} & \textbf{Qu\'e significa} & \textbf{Qu\'e debes saber defender} \\
\midrule
\endhead
\texttt{dni} / \(DNI\) & Irradiancia directa normal. Radiaci\'on solar directa sobre una superficie perpendicular al rayo solar. & Entra al c\'alculo de \(GHI\), pero no se usa como variable de entrada en XGBoost para evitar reconstruir indirectamente \(\fd\). \\
\texttt{dhi} / \(DHI\) & Irradiancia difusa horizontal. Radiaci\'on dispersada por atm\'osfera, nubes y aerosoles. & Es el numerador de la fracci\'on difusa. Es central para distinguir r\'egimen directo, mixto y difuso. \\
\texttt{ghi} / \(GHI\) & Irradiancia global horizontal. Suma de componente directa proyectada en horizontal y componente difusa. & En este trabajo se reconstruye como \(GHI = DNI\cos(\theta_z)+DHI\) usando pvlib y coordenadas efectivas de Santiago. \\
\texttt{poa} / \(POA\) & Irradiancia en el plano del arreglo fotovoltaico. & Es la irradiancia que realmente usa el panel inclinado. Es la entrada f\'isica principal del modelo NOCT. Viene directamente de PVWatts. \\
\texttt{temperature} & Temperatura ambiente horaria. & Se usa para calcular la temperatura de celda en NOCT. A mayor temperatura de celda, menor potencia por efecto t\'ermico. \\
\texttt{wind\_speed} & Velocidad del viento. & No entra en el NOCT simplificado, pero s\'i entra en XGBoost. Salir como variable importante sugiere que parte del residuo est\'a asociado a enfriamiento convectivo no modelado por NOCT. \\
\texttt{albedo} & Reflectancia del suelo. & PVWatts lo entrega; en esta etapa no es variable principal del modelo residual. \\
\texttt{zenith} & \'Angulo cenital solar. & Se usa para proyectar \(DNI\) sobre el plano horizontal al reconstruir \(GHI\). \\
\bottomrule
\end{longtable}

\subsection{Variable central: fracci\'on difusa}

\[
    \fd = \frac{DHI}{GHI}.
\]

La fracci\'on difusa mide qu\'e parte de la irradiancia global horizontal corresponde a radiaci\'on difusa. Si \(\fd\) es bajo, domina la radiaci\'on directa. Si \(\fd\) es alto, domina la componente difusa.

\begin{center}
\begin{tabular}{ll}
\toprule
\textbf{R\'egimen} & \textbf{Condici\'on} \\
\midrule
Baja fracci\'on difusa & \(\fd < 0{,}3\) \\
Mixto & \(0{,}3 \leq \fd < 0{,}7\) \\
Alta fracci\'on difusa & \(\fd \geq 0{,}7\) \\
\bottomrule
\end{tabular}
\end{center}

Defensa corta: \textbf{esta variable no se usa porque ``prediga todo'', sino porque permite diagnosticar si el error del modelo f\'isico cambia con el r\'egimen radiativo}. La tesis prueba si aporta informaci\'on adicional a un corrector residual.

\subsection{Variables de potencia y temperatura del modelo}

\begin{longtable}{p{3.2cm}p{5.6cm}p{5.4cm}}
\toprule
\textbf{Variable} & \textbf{Definici\'on} & \textbf{Uso en el proyecto} \\
\midrule
\endhead
\texttt{p\_ref\_dc} / \(\pref\) & Potencia DC entregada por PVWatts. & Referencia simulada. No es medici\'on real. Todo resultado debe defenderse bajo este alcance. \\
\texttt{tcell\_noct} / \(T_{c,\mathrm{NOCT}}\) & Temperatura de celda calculada con modelo NOCT. & Se calcula desde temperatura ambiente y \(POA\). \\
\texttt{p\_noct\_raw} & Potencia NOCT antes de p\'erdidas globales. & Variable auxiliar para revisar el modelo f\'isico sin p\'erdidas. \\
\texttt{p\_noct} / \(\pnoct\) & Potencia estimada por NOCT con p\'erdidas del 14\%. & Modelo f\'isico base que se compara contra PVWatts. \\
\texttt{epsilon} / \(\epsilon\) & Residuo definido como \(\epsilon = \pref-\pnoct\). & Variable objetivo del aprendizaje residual. Si \(\epsilon>0\), NOCT subestima PVWatts; si \(\epsilon<0\), NOCT sobreestima PVWatts. \\
\texttt{epsilon\_hat} / \(\widehat{\epsilon}\) & Residuo predicho por XGBoost. & Correcci\'on que se suma al modelo NOCT. \\
\texttt{p\_hybrid} / \(\phybrid\) & Potencia h\'ibrida corregida. & \(\phybrid = \pnoct + \widehat{\epsilon}\). \\
\bottomrule
\end{longtable}

Ecuaciones que debes dominar:

\[
    T_{c,\mathrm{NOCT}} = T_a + \left(\frac{NOCT-20}{800}\right)POA
\]

\[
    P_{\mathrm{NOCT,raw}} =
    P_{DC0}
    \left(\frac{POA}{G_{ref}}\right)
    \left[1+\gamma(T_{c,\mathrm{NOCT}}-T_{ref})\right]
\]

\[
    \pnoct = P_{\mathrm{NOCT,raw}}(1-L),
    \qquad
    \epsilon = \pref-\pnoct,
    \qquad
    \phybrid = \pnoct+\widehat{\epsilon}.
\]

Valores num\'ericos importantes:

\begin{itemize}
    \item \(P_{DC0} = 1000\ \mathrm{W}\), sistema de 1 kW DC.
    \item \(NOCT = 45^\circ C\).
    \item \(G_{ref} = 1000\ \mathrm{W/m^2}\).
    \item \(T_{ref} = 25^\circ C\).
    \item \(\gamma = -0{,}004\ ^\circ C^{-1}\).
    \item \(L = 0{,}14\), p\'erdidas globales del 14\%.
\end{itemize}

\subsection{Variables de filtrado}

\begin{longtable}{p{4.0cm}p{9.5cm}}
\toprule
\textbf{Filtro} & \textbf{Raz\'on} \\
\midrule
\endhead
\(GHI > 10\ \mathrm{W/m^2}\) & Evita calcular \(\fd = DHI/GHI\) con denominadores cercanos a cero. \\
\(POA > 10\ \mathrm{W/m^2}\) & Evita horas donde el panel casi no recibe irradiancia. \\
\(\pref > 0\) & Evita comparar horas sin generaci\'on DC simulada. \\
\(\pref > 50\ \mathrm{W}\) & Solo para error relativo. Evita porcentajes artificialmente grandes cuando la potencia es muy baja. \\
\texttt{is\_valid\_model} & Columna booleana que concentra los criterios para diagn\'ostico y modelado. \\
\bottomrule
\end{longtable}

\subsection{Variables de XGBoost}

Modelo base sin fracci\'on difusa:

\begin{itemize}
    \item \texttt{poa}
    \item \texttt{temperature}
    \item \texttt{wind\_speed}
    \item \texttt{hour\_sin}, \texttt{hour\_cos}
    \item \texttt{month\_sin}, \texttt{month\_cos}
\end{itemize}

Modelo con fracci\'on difusa:

\begin{itemize}
    \item Todas las variables anteriores.
    \item \texttt{fd}.
\end{itemize}

Punto clave de defensa: \textbf{el modelo sin \(\fd\) no usa \texttt{dni}, \texttt{dhi} ni \texttt{ghi}} porque desde esas variables se podr\'ia reconstruir \(\fd\). Si se incluyeran, la comparaci\'on con/sin fracci\'on difusa no ser\'ia limpia.

\section{Qu\'e hace exactamente este proyecto de t\'itulo}

\subsection{Idea central}

El proyecto eval\'ua si la fracci\'on difusa ayuda a explicar y corregir el error de un modelo f\'isico simplificado de potencia fotovoltaica.

La estructura real es:

\begin{enumerate}
    \item PVWatts genera una referencia simulada de potencia DC para Santiago.
    \item Se calcula un modelo f\'isico propio basado en NOCT.
    \item Se calcula el residuo entre PVWatts y NOCT.
    \item Se analiza si ese residuo cambia con la fracci\'on difusa.
    \item Se entrena XGBoost para aprender el residuo.
    \item Se compara un h\'ibrido sin \(\fd\) contra otro h\'ibrido con \(\fd\).
\end{enumerate}

La pregunta que debes defender no es ``si XGBoost es mejor que NOCT'', porque eso era esperable. La pregunta real es m\'as fina:

\[
    \text{Dado un corrector residual, \textquestiondown agregar } \fd
    \text{ mejora la correcci\'on?}
\]

\subsection{C\'odigo 1: descarga y preparaci\'on PVWatts}

Archivo: \texttt{01\_descargar\_pvwatts.ipynb}.

Este c\'odigo pide a PVWatts una simulaci\'on horaria para Santiago. Primero intenta con \texttt{dataset=nsrdb}; si falla, intenta con \texttt{dataset=intl}. En la ejecuci\'on final se usa \texttt{intl}, y la fuente clim\'atica efectiva corresponde a Santiago con distancia cero.

Fragmento clave de configuraci\'on:

\begin{lstlisting}[style=pythonstyle]
PVWATTS_URL = os.getenv(
    "PVWATTS_URL",
    "https://developer.nlr.gov/api/pvwatts/v8.json"
)

# Ubicacion: Santiago, Chile
LAT = -33.38000106811523
LON = -70.77999877929688

SYSTEM_CAPACITY = 1.0   # kW DC
AZIMUTH = 0             # norte, hemisferio sur
TILT = 20
ARRAY_TYPE = 1
MODULE_TYPE = 0
LOSSES = 14
\end{lstlisting}

La funci\'on de solicitud arma los par\'ametros para la API:

\begin{lstlisting}[style=pythonstyle]
def request_pvwatts(dataset: str) -> dict:
    params = {
        "api_key": API_KEY,
        "lat": LAT,
        "lon": LON,
        "system_capacity": SYSTEM_CAPACITY,
        "azimuth": AZIMUTH,
        "tilt": TILT,
        "array_type": ARRAY_TYPE,
        "module_type": MODULE_TYPE,
        "losses": LOSSES,
        "timeframe": "hourly",
        "dataset": dataset,
    }
\end{lstlisting}

El fallback metodol\'ogico est\'a aqu\'i:

\begin{lstlisting}[style=pythonstyle]
datasets_to_try = ["nsrdb", "intl"]

for dataset in datasets_to_try:
    try:
        data = request_pvwatts(dataset)
        data["_dataset_used"] = dataset
        return data
    except Exception as exc:
        print(f"Fallo con dataset='{dataset}': {exc}")
\end{lstlisting}

Defensa: no es que se haya mezclado Santiago con otra ciudad. Se usaron las coordenadas efectivas de Santiago y la fuente clim\'atica reportada por PVWatts queda a distancia cero.

Luego el c\'odigo reconstruye \(GHI\) y calcula \(\fd\):

\begin{lstlisting}[style=pythonstyle]
solar_position = pvlib.solarposition.get_solarposition(
    time=df.index,
    latitude=SOLAR_GEOMETRY_LAT,
    longitude=SOLAR_GEOMETRY_LON,
)

cos_zenith = np.cos(np.deg2rad(solar_position["zenith"]))
cos_zenith = np.clip(cos_zenith, 0, None)

df["ghi"] = df["dni"] * cos_zenith + df["dhi"]
df["ghi"] = df["ghi"].clip(lower=0)

df["fd"] = np.nan
mask = df["ghi"] > GHI_MIN_FOR_FD
df.loc[mask, "fd"] = df.loc[mask, "dhi"] / df.loc[mask, "ghi"]
df.loc[(df["fd"] < 0) | (df["fd"] > 1), "fd"] = np.nan
\end{lstlisting}

Clasificaci\'on de r\'egimenes:

\begin{lstlisting}[style=pythonstyle]
df["regimen_fd"] = pd.cut(
    df["fd"],
    bins=[-np.inf, 0.3, 0.7, np.inf],
    labels=["baja_fd", "mixto", "alta_fd"],
    right=False,
)
\end{lstlisting}

Salida:

\begin{itemize}
    \item \texttt{data/01\_pvwatts\_santiago.csv}
    \item \texttt{data/01\_pvwatts\_santiago\_metadata.json}
\end{itemize}

Resultado base:

\begin{itemize}
    \item 8760 filas.
    \item 365 d\'ias.
    \item 17 columnas principales.
    \item Dataset final: \texttt{intl}.
\end{itemize}

\subsection{C\'odigo 2: modelo NOCT y residuo}

Archivo: \texttt{02\_noct\_vs\_pvwatts.ipynb}.

Este c\'odigo toma el CSV del c\'odigo 1, calcula la potencia NOCT y define el residuo. Aqu\'i nace la variable objetivo de XGBoost.

Par\'ametros:

\begin{lstlisting}[style=pythonstyle]
PDC0_W = 1000.0
SYSTEM_LOSSES = 0.14

G_REF = 1000.0
T_REF = 25.0
NOCT = 45.0
GAMMA = -0.004
\end{lstlisting}

Modelo NOCT:

\begin{lstlisting}[style=pythonstyle]
df["tcell_noct"] = (
    df["temperature"] + ((NOCT - 20.0) / 800.0) * df["poa"]
)

temperature_factor = 1 + GAMMA * (df["tcell_noct"] - T_REF)

df["p_noct_raw"] = (
    PDC0_W
    * (df["poa"] / G_REF)
    * temperature_factor
)

df["p_noct"] = df["p_noct_raw"] * (1 - SYSTEM_LOSSES)
\end{lstlisting}

Interpretaci\'on:

\begin{itemize}
    \item \(POA\) aumenta potencia.
    \item Temperatura de celda alta reduce potencia porque \(\gamma<0\).
    \item Se aplican p\'erdidas del 14\% para hacer NOCT comparable con la salida DC de PVWatts.
\end{itemize}

Definici\'on de referencia y residuo:

\begin{lstlisting}[style=pythonstyle]
df["p_ref_dc"] = df["dc"]
df["epsilon"] = df["p_ref_dc"] - df["p_noct"]
df["abs_epsilon"] = df["epsilon"].abs()
\end{lstlisting}

Lectura del signo de \(\epsilon\):

\begin{itemize}
    \item \(\epsilon>0\): PVWatts es mayor que NOCT. Entonces NOCT subestima.
    \item \(\epsilon<0\): PVWatts es menor que NOCT. Entonces NOCT sobreestima.
\end{itemize}

C\'alculo de error relativo:

\begin{lstlisting}[style=pythonstyle]
df["relative_error"] = np.nan
mask_relative = df["p_ref_dc"] > P_REF_MIN_RELATIVE_ERROR

df.loc[mask_relative, "relative_error"] = (
    df.loc[mask_relative, "abs_epsilon"]
    / df.loc[mask_relative, "p_ref_dc"]
)
\end{lstlisting}

Filtro de observaciones v\'alidas:

\begin{lstlisting}[style=pythonstyle]
df["is_valid_model"] = (
    (df["ghi"] > GHI_MIN)
    & (df["poa"] > POA_MIN)
    & (df["p_ref_dc"] > 0)
    & df["fd"].notna()
    & df["regimen_fd"].notna()
)
\end{lstlisting}

Salida:

\begin{itemize}
    \item \texttt{data/02\_noct\_vs\_pvwatts\_santiago.csv}
    \item \texttt{data/02\_noct\_vs\_pvwatts\_santiago\_metadata.json}
\end{itemize}

Resultado:

\begin{itemize}
    \item 8760 filas totales.
    \item 4215 filas v\'alidas para diagn\'ostico/modelado.
    \item 3688 filas con error relativo definido.
\end{itemize}

\subsection{C\'odigo 3: diagn\'ostico del residuo}

Archivo: \texttt{03\_diagnostico\_fd.ipynb}.

Este c\'odigo no entrena modelos. Su trabajo es responder: \textbf{\textquestiondown el error del NOCT cambia con la fracci\'on difusa?}

Lee la base del c\'odigo 2:

\begin{lstlisting}[style=pythonstyle]
INPUT_FILE = Path("data/02_noct_vs_pvwatts_santiago.csv")
OUT_DIR = Path("outputs/03_diagnostico_fd_santiago")
\end{lstlisting}

Funci\'on de m\'etricas:

\begin{lstlisting}[style=pythonstyle]
def regression_metrics(df):
    y_true = df["p_ref_dc"]
    y_pred = df["p_noct"]
    error = y_pred - y_true

    return {
        "n": len(df),
        "MAE_W": mean_absolute_error(y_true, y_pred),
        "RMSE_W": np.sqrt(mean_squared_error(y_true, y_pred)),
        "Bias_W": error.mean(),
        "R2": r2_score(y_true, y_pred),
        "Error_relativo_promedio": df["relative_error"].mean(),
        "fd_promedio": df["fd"].mean(),
    }
\end{lstlisting}

Ojo con el bias: en diagn\'ostico se reporta como \(P_{modelo}-P_{ref}\). Por eso un bias negativo significa que el modelo subestima la referencia.

Correlaciones:

\begin{lstlisting}[style=pythonstyle]
variables_error = ["epsilon", "abs_epsilon", "relative_error"]

for var in variables_error:
    pearson = df_valid["fd"].corr(df_valid[var], method="pearson")
    spearman = df_valid["fd"].corr(df_valid[var], method="spearman")
\end{lstlisting}

Gr\'aficos generados:

\begin{itemize}
    \item \texttt{scatter\_epsilon\_vs\_fd\_santiago.png}
    \item \texttt{boxplot\_abs\_error\_por\_regimen\_santiago.png}
    \item \texttt{boxplot\_error\_relativo\_por\_regimen\_santiago.png}
    \item \texttt{scatter\_pnoct\_vs\_pref\_santiago.png}
    \item \texttt{bar\_mae\_por\_regimen\_santiago.png}
\end{itemize}

Resultados que debes memorizar:

\begin{center}
\begin{tabular}{lr}
\toprule
\textbf{M\'etrica global NOCT} & \textbf{Valor} \\
\midrule
MAE & 11,05 W \\
RMSE & 15,05 W \\
Bias \(P_{NOCT}-P_{ref}\) & -1,60 W \\
\(R^2\) & 0,9968 \\
Error relativo promedio & 3,72\% \\
\(\fd\) promedio & 0,5797 \\
\bottomrule
\end{tabular}
\end{center}

Por r\'egimen:

\begin{center}
\begin{tabular}{lrrrr}
\toprule
\textbf{R\'egimen} & \textbf{n} & \textbf{MAE [W]} & \textbf{Bias [W]} & \textbf{Error rel. [\%]} \\
\midrule
Baja \(\fd\) & 1349 & 20,42 & -14,54 & 2,98 \\
Mixto & 1126 & 7,88 & 2,51 & 3,07 \\
Alta \(\fd\) & 1740 & 5,83 & 5,76 & 5,15 \\
\bottomrule
\end{tabular}
\end{center}

Lectura correcta:

\begin{itemize}
    \item En baja \(\fd\), el error absoluto es mayor porque hay m\'as potencia disponible. Errores proporcionales peque\~nos se vuelven varios watts.
    \item En alta \(\fd\), el error absoluto baja, pero el error relativo sube porque la potencia disponible es menor.
    \item El cambio de signo del bias indica estructura del residuo, no solo cambio de escala.
\end{itemize}

Correlaciones:

\begin{center}
\begin{tabular}{lrr}
\toprule
\textbf{Variable} & \textbf{Pearson con \(\fd\)} & \textbf{Spearman con \(\fd\)} \\
\midrule
\(\epsilon\) & -0,5610 & -0,5015 \\
\(|\epsilon|\) & -0,5776 & -0,5729 \\
Error relativo & 0,3591 & 0,3119 \\
\bottomrule
\end{tabular}
\end{center}

Defensa: estas correlaciones no prueban causalidad, pero s\'i muestran asociaci\'on sistem\'atica entre fracci\'on difusa y error.

\subsection{C\'odigo 4: modelo h\'ibrido con XGBoost}

Archivo: \texttt{04\_modelo\_hibrido\_xgboost.ipynb}.

Este es el coraz\'on de la tesis. Toma la base con residuo y entrena dos modelos XGBoost:

\begin{itemize}
    \item H\'ibrido sin fracci\'on difusa.
    \item H\'ibrido con fracci\'on difusa.
\end{itemize}

Configuraci\'on principal:

\begin{lstlisting}[style=pythonstyle]
INPUT_FILE = Path("data/02_noct_vs_pvwatts_santiago.csv")
OUT_DIR = Path("outputs/04_modelo_hibrido_xgboost_santiago")

TARGET = "epsilon"

FEATURES_BASE = [
    "poa",
    "temperature",
    "wind_speed",
    "hour_sin",
    "hour_cos",
    "month_sin",
    "month_cos",
]

FEATURES_WITH_FD = FEATURES_BASE + ["fd"]
\end{lstlisting}

Clave: \texttt{TARGET = "epsilon"}. XGBoost no aprende directamente la potencia, aprende el error del NOCT.

Hiperpar\'ametros:

\begin{lstlisting}[style=pythonstyle]
XGB_PARAMS = {
    "objective": "reg:squarederror",
    "n_estimators": 350,
    "max_depth": 3,
    "learning_rate": 0.035,
    "subsample": 0.85,
    "colsample_bytree": 0.85,
    "min_child_weight": 3,
    "reg_lambda": 2.0,
    "random_state": 42,
    "n_jobs": -1,
}
\end{lstlisting}

Transformaci\'on temporal:

\begin{lstlisting}[style=pythonstyle]
model_df["hour_sin"] = np.sin(2 * np.pi * model_df["hour"] / 24)
model_df["hour_cos"] = np.cos(2 * np.pi * model_df["hour"] / 24)

model_df["month_sin"] = np.sin(2 * np.pi * (model_df["month"] - 1) / 12)
model_df["month_cos"] = np.cos(2 * np.pi * (model_df["month"] - 1) / 12)
\end{lstlisting}

Por qu\'e se hace: hora 23 y hora 0 est\'an cerca en el ciclo diario. Si usaras la hora como n\'umero lineal, 23 y 0 parecer\'ian extremos opuestos.

Separaci\'on temporal:

\begin{lstlisting}[style=pythonstyle]
split_index = int(len(model_df) * (1 - TEST_FRACTION))
train_df = model_df.iloc[:split_index].copy()
test_df = model_df.iloc[split_index:].copy()
\end{lstlisting}

No se usa split aleatorio porque los datos horarios cercanos se parecen entre s\'i. Si mezclas horas vecinas entre entrenamiento y prueba, el resultado puede verse artificialmente bueno.

Entrenamiento:

\begin{lstlisting}[style=pythonstyle]
def fit_xgb_residual(train, features):
    model = XGBRegressor(**XGB_PARAMS)
    model.fit(train[features], train[TARGET])
    return model

model_sin_fd = fit_xgb_residual(train_df, FEATURES_BASE)
model_con_fd = fit_xgb_residual(train_df, FEATURES_WITH_FD)
\end{lstlisting}

Predicci\'on residual y potencia corregida:

\begin{lstlisting}[style=pythonstyle]
def add_predictions(df_eval, model, features, suffix):
    df_out = df_eval.copy()
    eps_hat_col = f"epsilon_hat_{suffix}"
    p_hybrid_col = f"p_hybrid_{suffix}"

    df_out[eps_hat_col] = model.predict(df_out[features])
    df_out[p_hybrid_col] = df_out["p_noct"] + df_out[eps_hat_col]
    df_out[p_hybrid_col] = df_out[p_hybrid_col].clip(lower=0)
    return df_out[[eps_hat_col, p_hybrid_col]]
\end{lstlisting}

Evaluaci\'on:

\begin{lstlisting}[style=pythonstyle]
def evaluate_models(df_eval, dataset_name):
    rows = []
    rows.append(evaluate_power(df_eval["p_ref_dc"], df_eval["p_noct"], MODEL_NOCT, dataset_name))
    rows.append(evaluate_power(df_eval["p_ref_dc"], df_eval["p_hybrid_sin_fd"], MODEL_XGB_SIN_FD, dataset_name))
    rows.append(evaluate_power(df_eval["p_ref_dc"], df_eval["p_hybrid_con_fd"], MODEL_XGB_CON_FD, dataset_name))

    metrics = pd.DataFrame(rows)
    noct_mae = metrics.loc[metrics["modelo"] == MODEL_NOCT, "MAE_W"].iloc[0]
    metrics["Mejora_MAE_vs_NOCT_pct"] = (
        (noct_mae - metrics["MAE_W"]) / noct_mae * 100
    )
    return metrics
\end{lstlisting}

Validaci\'on temporal:

\begin{lstlisting}[style=pythonstyle]
def run_time_series_cv(df_cv, features, model_name, n_splits=5):
    tscv = TimeSeriesSplit(n_splits=n_splits)
    rows = []
    for fold, (train_idx, test_idx) in enumerate(tscv.split(df_cv), start=1):
        fold_train = df_cv.iloc[train_idx].copy()
        fold_test = df_cv.iloc[test_idx].copy()
        model = fit_xgb_residual(fold_train, features)
        pred = add_predictions(fold_test, model, features, "cv")
        fold_eval = pd.concat([fold_test, pred], axis=1)
        rows.append(evaluate_power(
            fold_eval["p_ref_dc"],
            fold_eval["p_hybrid_cv"],
            model_name,
            f"cv_fold_{fold}",
        ))
    return pd.DataFrame(rows)
\end{lstlisting}

Importancia de variables:

\begin{lstlisting}[style=pythonstyle]
importance_con_fd = pd.DataFrame({
    "feature": FEATURES_WITH_FD,
    "importance": model_con_fd.feature_importances_,
}).sort_values("importance", ascending=False)
\end{lstlisting}

Resultados de test:

\begin{center}
\begin{tabular}{lrrr}
\toprule
\textbf{Modelo} & \textbf{MAE [W]} & \textbf{RMSE [W]} & \textbf{Mejora MAE vs NOCT} \\
\midrule
NOCT & 13,30 & 17,10 & 0,00\% \\
H\'ibrido sin \(\fd\) & 2,70 & 3,73 & 79,67\% \\
H\'ibrido con \(\fd\) & 2,22 & 3,01 & 83,27\% \\
\bottomrule
\end{tabular}
\end{center}

Lectura correcta:

\begin{itemize}
    \item XGBoost sin \(\fd\) ya mejora mucho porque aprende estructura residual usando \(POA\), temperatura, viento y ciclos temporales.
    \item La comparaci\'on importante es sin \(\fd\) versus con \(\fd\).
    \item Pasar de 2,70 W a 2,22 W equivale a una mejora adicional aproximada de 17,7\% respecto al h\'ibrido sin \(\fd\).
\end{itemize}

Por r\'egimen en test:

\begin{center}
\begin{tabular}{lrrr}
\toprule
\textbf{R\'egimen} & \textbf{NOCT [W]} & \textbf{H\'ibrido con \(\fd\) [W]} & \textbf{Mejora} \\
\midrule
Baja \(\fd\) & 21,13 & 2,03 & 90,40\% \\
Mixto & 9,27 & 2,95 & 68,14\% \\
Alta \(\fd\) & 6,16 & 1,83 & 70,23\% \\
\bottomrule
\end{tabular}
\end{center}

Validaci\'on temporal:

\begin{center}
\begin{tabular}{lrr}
\toprule
\textbf{Modelo} & \textbf{MAE promedio [W]} & \textbf{RMSE promedio [W]} \\
\midrule
H\'ibrido sin \(\fd\) & \(3,37 \pm 1,32\) & \(4,44 \pm 1,60\) \\
H\'ibrido con \(\fd\) & \(2,52 \pm 0,64\) & \(3,36 \pm 0,85\) \\
\bottomrule
\end{tabular}
\end{center}

Adem\'as, el modelo con \(\fd\) gana en los 5 pliegues temporales. Esa frase es importante para defensa.

Importancia de variables del modelo con \(\fd\):

\begin{center}
\begin{tabular}{lr}
\toprule
\textbf{Variable} & \textbf{Importancia} \\
\midrule
Fracci\'on difusa & 0,299 \\
Velocidad de viento & 0,239 \\
\(POA\) & 0,157 \\
Temperatura ambiente & 0,142 \\
Hora seno & 0,062 \\
\bottomrule
\end{tabular}
\end{center}

Interpretaci\'on:

\begin{itemize}
    \item \(\fd\) aparece como la variable m\'as importante del modelo con fracci\'on difusa.
    \item Viento aparece segundo, lo que tiene sentido f\'isico porque NOCT no incluye enfriamiento convectivo expl\'icito.
    \item Importancia de XGBoost no es causalidad; es diagn\'ostico predictivo interno del ensamble.
\end{itemize}

\section{Resumen mental para defender el proyecto}

Si te piden explicar la tesis en un minuto:

\begin{quote}
Mi proyecto eval\'ua si la fracci\'on difusa de la irradiancia solar ayuda a diagnosticar y corregir el error de un modelo f\'isico NOCT para potencia fotovoltaica. Primero descargu\'e una simulaci\'on horaria de PVWatts para Santiago. Luego calcul\'e un modelo NOCT propio con \(POA\), temperatura ambiente, coeficiente t\'ermico y p\'erdidas globales. Defin\'i el residuo como \(P_{ref}-P_{NOCT}\), donde \(P_{ref}\) es la potencia DC simulada por PVWatts. Despu\'es analic\'e c\'omo cambia ese residuo seg\'un la fracci\'on difusa \(DHI/GHI\). Finalmente entren\'e dos modelos XGBoost para aprender el residuo: uno sin fracci\'on difusa y otro con fracci\'on difusa. El modelo con fracci\'on difusa redujo el MAE de prueba desde 2,70 W a 2,22 W respecto al h\'ibrido sin fracci\'on difusa, y gan\'o en los cinco pliegues de validaci\'on temporal. Por eso concluyo que, dentro del marco simulado de PVWatts, la fracci\'on difusa aporta informaci\'on adicional para corregir el modelo f\'isico.
\end{quote}

\section{Preguntas realistas de comit\'e y respuestas}

\subsection{Sobre el objetivo y la hip\'otesis}

\textbf{1. \textquestiondown Cu\'al es exactamente la pregunta de investigaci\'on?}

La pregunta es si la fracci\'on difusa \(\fd=DHI/GHI\) se relaciona con el error residual de un modelo f\'isico NOCT y si incorporarla en un modelo residual supervisado mejora la estimaci\'on de potencia fotovoltaica respecto al modelo f\'isico base.

\textbf{2. \textquestiondown Cu\'al fue la hip\'otesis?}

Que el residuo del NOCT no es completamente aleatorio, sino que presenta estructura asociada al r\'egimen de fracci\'on difusa. Si eso es cierto, \(\fd\) deber\'ia ayudar tanto al diagn\'ostico del error como a la correcci\'on residual con XGBoost.

\textbf{3. \textquestiondown La hip\'otesis se confirma?}

S\'i, dentro del alcance simulado de PVWatts. Hay tres evidencias: el residuo cambia con \(\fd\), las correlaciones entre \(\fd\) y error son no nulas, y el modelo h\'ibrido con \(\fd\) mejora al h\'ibrido sin \(\fd\) tanto en el holdout temporal como en los cinco pliegues de validaci\'on temporal.

\textbf{4. \textquestiondown Qu\'e ser\'ia un resultado contrario a tu hip\'otesis?}

Que el residuo no cambiara con \(\fd\), que las correlaciones fueran cercanas a cero, o que el modelo con \(\fd\) no mejorara al modelo sin \(\fd\). En ese caso habr\'ia concluido que la fracci\'on difusa no aporta informaci\'on adicional bajo esta configuraci\'on.

\subsection{Sobre PVWatts y los datos}

\textbf{5. \textquestiondown PVWatts es una medici\'on real?}

No. PVWatts entrega una simulaci\'on. En la tesis se usa como referencia simulada, no como verdad experimental de campo. Eso est\'a declarado como alcance y limitaci\'on.

\textbf{6. \textquestiondown Entonces qu\'e demuestras realmente?}

Demuestro una propiedad metodol\'ogica: que el error de un modelo f\'isico simplificado frente a una referencia simulada m\'as completa tiene estructura asociada a la fracci\'on difusa, y que esa estructura puede ser aprendida por un corrector residual. No afirmo validaci\'on experimental con datos reales.

\textbf{7. \textquestiondown Por qu\'e usaste Santiago?}

Porque se decidi\'o trabajar con una ubicaci\'on consistente espacialmente. PVWatts entrega datos para Santiago con distancia cero respecto a las coordenadas efectivas usadas, y esas mismas coordenadas se usan para reconstruir \(GHI\) con pvlib.

\textbf{8. \textquestiondown Qu\'e pas\'o con NSRDB?}

La API no encontr\'o datos con \texttt{dataset=nsrdb} para esas coordenadas. Por eso el c\'odigo tiene un fallback hacia \texttt{dataset=intl}. El dataset final usado fue \texttt{intl}.

\textbf{9. \textquestiondown Usar \texttt{intl} debilita el trabajo?}

No invalida el trabajo, pero delimita su alcance. El an\'alisis sigue siendo reproducible y consistente espacialmente, pero no se debe presentar como una validaci\'on experimental local.

\textbf{10. \textquestiondown Por qu\'e reconstruyes \(GHI\) si PVWatts entrega otras variables?}

Porque la fracci\'on difusa requiere \(GHI\). PVWatts entrega \(DNI\), \(DHI\) y \(POA\); entonces \(GHI\) se reconstruye con \(GHI=DNI\cos(\theta_z)+DHI\), usando la posici\'on solar calculada con pvlib.

\textbf{11. \textquestiondown Qu\'e significa que el a\~no sea meteorol\'ogico t\'ipico?}

Significa que las 8760 horas representan un a\~no t\'ipico construido a partir de datos clim\'aticos, no una serie hist\'orica real de un a\~no espec\'ifico. Las fechas se usan como \'indice temporal para ordenar la serie.

\subsection{Sobre el modelo NOCT}

\textbf{12. \textquestiondown Por qu\'e usaste NOCT?}

Porque es un modelo f\'isico simple, interpretable y adecuado como baseline. Permite estimar temperatura de celda y potencia usando \(POA\), temperatura ambiente y par\'ametros del m\'odulo.

\textbf{13. \textquestiondown Qu\'e limita al NOCT?}

Es simplificado. No representa expl\'icitamente la composici\'on directa-difusa, ni incluye un t\'ermino expl\'icito de viento/convecci\'on en la versi\'on usada. Por eso puede dejar un residuo estructurado.

\textbf{14. \textquestiondown Qu\'e significa \(\gamma=-0{,}004\)?}

Es el coeficiente t\'ermico del m\'odulo. Como es negativo, cuando la temperatura de celda sube por encima de \(25^\circ C\), la potencia disminuye. El valor \(-0{,}004\ ^\circ C^{-1}\) equivale aproximadamente a una p\'erdida de 0,4\% por grado Celsius.

\textbf{15. \textquestiondown Por qu\'e aplicas p\'erdidas del 14\% al NOCT?}

Porque PVWatts fue configurado con p\'erdidas globales del 14\%. Aplicarlas tambi\'en al NOCT hace que ambas potencias DC sean comparables bajo el mismo supuesto de p\'erdidas.

\textbf{16. \textquestiondown Por qu\'e comparas contra DC y no AC?}

Porque el modelo NOCT implementado estima potencia DC del arreglo. Compararlo contra AC introducir\'ia efectos de inversor que no forman parte del modelo f\'isico usado.

\textbf{17. \textquestiondown Tu NOCT replica PVWatts?}

No exactamente. El NOCT es un modelo f\'isico simplificado propio. PVWatts tiene submodelos internos m\'as completos. La diferencia entre ambos es precisamente el residuo que se analiza.

\subsection{Sobre fracci\'on difusa}

\textbf{18. \textquestiondown Por qu\'e te importa la fracci\'on difusa?}

Porque dos horas pueden tener irradiancia sobre el plano similar, pero distinta composici\'on directa-difusa. La distribuci\'on angular y atmosf\'erica de la radiaci\'on puede cambiar el comportamiento del modelo y su error.

\textbf{19. \textquestiondown Por qu\'e los umbrales 0,3 y 0,7?}

Se usan como una partici\'on operativa simple: baja fracci\'on difusa, r\'egimen mixto y alta fracci\'on difusa. No se presentan como fronteras universales, sino como una clasificaci\'on reproducible para estudiar el comportamiento del residuo.

\textbf{20. \textquestiondown Por qu\'e el error absoluto es mayor en baja \(\fd\), pero el relativo mayor en alta \(\fd\)?}

Porque miden cosas distintas. El error absoluto mide watts. En baja \(\fd\) suele haber m\'as irradiancia y potencia, entonces una diferencia proporcional peque\~na puede ser muchos watts. El error relativo divide por la potencia de referencia. En alta \(\fd\), la potencia suele ser menor, entonces pocos watts pueden representar un porcentaje mayor.

\textbf{21. \textquestiondown El cambio de bias por r\'egimen qu\'e indica?}

Indica que el error no solo cambia de magnitud, sino tambi\'en de direcci\'on promedio. En baja \(\fd\), NOCT tiende a subestimar; en alta \(\fd\), tiende a sobreestimar. Eso respalda que el residuo tiene estructura.

\textbf{22. \textquestiondown La correlaci\'on prueba causalidad?}

No. La correlaci\'on muestra asociaci\'on estad\'istica. La causalidad f\'isica requerir\'ia un dise\~no experimental o mediciones adicionales. En esta tesis se usa como evidencia diagn\'ostica, no causal.

\subsection{Sobre XGBoost y aprendizaje residual}

\textbf{23. \textquestiondown Por qu\'e XGBoost y no una red neuronal?}

Porque la base es tabular, de tama\~no moderado, y con variables f\'isicas estructuradas. XGBoost suele funcionar muy bien en datos tabulares, requiere menos datos que una red profunda y permite inspeccionar importancias de variables.

\textbf{24. \textquestiondown Esto es aprendizaje profundo?}

No. Es aprendizaje autom\'atico supervisado con XGBoost, un ensamble de \'arboles de decisi\'on entrenados por gradient boosting. Aprendizaje profundo se refiere a redes neuronales con m\'ultiples capas, que no se usaron aqu\'i.

\textbf{25. \textquestiondown Qu\'e aprende XGBoost exactamente?}

Aprende \(\epsilon=P_{ref}-P_{NOCT}\), es decir, el residuo del modelo f\'isico. Luego la potencia corregida se calcula como \(P_{hybrid}=P_{NOCT}+\widehat{\epsilon}\).

\textbf{26. \textquestiondown Por qu\'e aprender el residuo y no la potencia directamente?}

Porque el NOCT ya captura la relaci\'on f\'isica dominante entre irradiancia, temperatura y potencia. El modelo de datos se enfoca solo en corregir lo que el modelo f\'isico no representa bien. Eso mantiene interpretabilidad y reduce la carga del aprendizaje supervisado.

\textbf{27. \textquestiondown C\'omo evitaste fuga de informaci\'on?}

El modelo sin \(\fd\) no recibe \(DNI\), \(DHI\) ni \(GHI\), porque desde ellas podr\'ia reconstruir \(\fd\). Adem\'as, la evaluaci\'on usa separaci\'on temporal, no mezcla aleatoria de horas cercanas.

\textbf{28. \textquestiondown Usar \(\fd\) contempor\'aneo es fuga?}

No para una tarea de estimaci\'on. \(\fd\) se usa como variable contempor\'anea de la misma hora, no como informaci\'on futura. Si el objetivo fuera pron\'ostico operacional, habr\'ia que usar pron\'osticos meteorol\'ogicos de \(\fd\) o de sus componentes.

\textbf{29. \textquestiondown Por qu\'e transformaste hora y mes con seno/coseno?}

Porque hora y mes son variables c\'iclicas. La hora 23 y la hora 0 son cercanas, pero como n\'umeros lineales parecer\'ian lejanas. Seno y coseno conservan esa continuidad circular.

\textbf{30. \textquestiondown Qu\'e hiperpar\'ametros usaste y por qu\'e?}

Us\'e una configuraci\'on moderada y regularizada: 350 \'arboles, profundidad m\'axima 3, learning rate 0,035, \texttt{subsample}=0,85, \texttt{colsample\_bytree}=0,85, \texttt{min\_child\_weight}=3 y \(\lambda=2\). No busqu\'e optimizar al m\'aximo XGBoost; la idea era comparar limpiamente dos modelos equivalentes salvo por \(\fd\).

\subsection{Sobre resultados}

\textbf{31. \textquestiondown Cu\'al es el resultado principal?}

El h\'ibrido con fracci\'on difusa obtiene menor error que el h\'ibrido sin fracci\'on difusa. En test, el MAE baja de 2,70 W a 2,22 W. Eso es una mejora adicional de aproximadamente 17,7\% respecto al h\'ibrido sin \(\fd\).

\textbf{32. \textquestiondown El h\'ibrido sin \(\fd\) ya mejora mucho. Entonces \(\fd\) importa poco?}

No. Justamente esa comparaci\'on hace m\'as exigente la prueba. El modelo sin \(\fd\) ya tiene variables fuertes como \(POA\), temperatura y viento. Aun as\'i, agregar \(\fd\) mejora el resultado, lo que indica aporte incremental.

\textbf{33. \textquestiondown Por qu\'e el \(R^2\) ya era alto con NOCT?}

Porque tanto NOCT como PVWatts dependen fuertemente de la irradiancia. Un \(R^2\) alto indica que la tendencia global se captura bien, pero no descarta errores sistem\'aticos por r\'egimen radiativo.

\textbf{34. \textquestiondown Qu\'e significa que el modelo con \(\fd\) gane en los cinco pliegues temporales?}

Significa que la mejora no aparece solo en un bloque particular de prueba. Al evaluar cinco particiones temporales crecientes, el modelo con \(\fd\) mantiene menor MAE en todas, lo que refuerza la estabilidad del aporte.

\textbf{35. \textquestiondown Qu\'e variable fue m\'as importante en XGBoost?}

En el modelo con fracci\'on difusa, \(\fd\) fue la variable m\'as importante, seguida por velocidad de viento, \(POA\) y temperatura ambiente.

\textbf{36. \textquestiondown Por qu\'e viento sale tan importante?}

Porque el viento afecta la temperatura de operaci\'on del m\'odulo por enfriamiento convectivo. El NOCT simplificado usado no incluye expl\'icitamente viento, entonces XGBoost puede usarlo para corregir diferencias t\'ermicas no modeladas.

\textbf{37. \textquestiondown Las importancias de XGBoost prueban mecanismos f\'isicos?}

No. Son indicadores predictivos internos del modelo. Sirven para diagnosticar qu\'e variables usa el ensamble, pero no deben confundirse con causalidad.

\subsection{Sobre validaci\'on y limitaciones}

\textbf{38. \textquestiondown Tu validaci\'on es suficiente?}

Es suficiente para el objetivo metodol\'ogico de esta etapa: comparar modelos bajo una referencia simulada y una separaci\'on temporal. No es suficiente para afirmar desempe\~no experimental real; eso queda como trabajo futuro.

\textbf{39. \textquestiondown Qu\'e falta para llevarlo a una aplicaci\'on real?}

Validar con mediciones reales de potencia, irradiancia y meteorolog\'ia. Adem\'as, para pron\'ostico operacional habr\'ia que usar predicciones meteorol\'ogicas futuras y evaluar incertidumbre.

\textbf{40. \textquestiondown El modelo sirve para pronosticar ma\~nana?}

No directamente. Este trabajo hace estimaci\'on/correcci\'on residual con variables contempor\'aneas. Para pronosticar ma\~nana se necesitar\'ian pron\'osticos de irradiancia, temperatura, viento y posiblemente fracci\'on difusa.

\textbf{41. \textquestiondown Por qu\'e no usaste mediciones reales?}

Porque esta etapa se plante\'o como un estudio metodol\'ogico reproducible usando PVWatts como referencia simulada. Las mediciones reales son el paso natural para validar externamente la metodolog\'ia.

\textbf{42. \textquestiondown Qu\'e pasa si PVWatts y NOCT tienen supuestos parecidos?}

Eso puede explicar el alto ajuste global. Sin embargo, la tesis no se basa solo en que NOCT se parezca a PVWatts, sino en el residuo restante y en si \(\fd\) ayuda a corregirlo.

\textbf{43. \textquestiondown Tus resultados son generalizables a otras ciudades?}

No autom\'aticamente. El estudio usa Santiago, una configuraci\'on fija y un a\~no meteorol\'ogico t\'ipico. Para generalizar se deben repetir pruebas en otras ubicaciones y configuraciones.

\textbf{44. \textquestiondown Por qu\'e no calculaste \(POA\) independientemente?}

En esta etapa se us\'o \(POA\) de PVWatts para aislar el diagn\'ostico del residuo NOCT frente a la referencia simulada. Calcular \(POA\) independientemente ser\'ia una extensi\'on \'util para separar errores de geometr\'ia, transposici\'on e irradiancia.

\textbf{45. \textquestiondown Cu\'al es la principal debilidad de la tesis?}

La principal limitaci\'on es que la referencia es simulada, no experimental. Por eso los resultados deben interpretarse como evidencia metodol\'ogica frente a PVWatts y no como validaci\'on de campo.

\subsection{Preguntas sobre c\'odigo}

\textbf{46. \textquestiondown Qu\'e archivo genera la base inicial?}

\texttt{01\_descargar\_pvwatts.ipynb}. Genera \texttt{data/01\_pvwatts\_santiago.csv} y su metadata.

\textbf{47. \textquestiondown Qu\'e archivo define el residuo?}

\texttt{02\_noct\_vs\_pvwatts.ipynb}. Define \texttt{epsilon = p\_ref\_dc - p\_noct}.

\textbf{48. \textquestiondown Qu\'e archivo genera los gr\'aficos de diagn\'ostico?}

\texttt{03\_diagnostico\_fd.ipynb}. Genera m\'etricas globales, m\'etricas por r\'egimen, correlaciones y gr\'aficos del error.

\textbf{49. \textquestiondown Qu\'e archivo entrena XGBoost?}

\texttt{04\_modelo\_hibrido\_xgboost.ipynb}. Entrena los modelos residuales sin y con fracci\'on difusa, eval\'ua test, validaci\'on temporal, m\'etricas por r\'egimen e importancia de variables.

\textbf{50. \textquestiondown Si tuvieras que revisar un solo c\'odigo para defender la tesis, cu\'al ser\'ia?}

El c\'odigo 4, porque contiene el modelo h\'ibrido y la comparaci\'on central. Pero para entenderlo bien hay que tener claro que el c\'odigo 2 define \(\epsilon\), que es la variable objetivo.

\section{Frases cortas que conviene memorizar}

\begin{itemize}
    \item ``PVWatts se usa como referencia simulada, no como medici\'on experimental.''
    \item ``El modelo h\'ibrido aprende el residuo del NOCT, no la potencia completa desde cero.''
    \item ``La comparaci\'on clave es h\'ibrido sin fracci\'on difusa versus h\'ibrido con fracci\'on difusa.''
    \item ``No inclu\'i \(DNI\), \(DHI\) ni \(GHI\) en el modelo sin \(\fd\) para evitar reconstruir indirectamente la fracci\'on difusa.''
    \item ``El error absoluto y el relativo cuentan historias distintas porque dependen de la escala de potencia.''
    \item ``La importancia de variables de XGBoost es diagn\'ostica, no causal.''
    \item ``El viento aparece como variable relevante porque el NOCT simplificado no incluye enfriamiento convectivo expl\'icito.''
    \item ``El resultado est\'a respaldado por el holdout temporal y por validaci\'on temporal en cinco particiones.''
\end{itemize}

\end{document}
